\originalchapter{自由Fermi气体与电子}\label{ch:fermi}
\originalsection{态密度与Fermi面}
一盒自由粒子的单粒子能量$\epsilon=\hbar^2k^2/(2m)$. 周期边界允许$k_i=2\pi n_i/L_i$, 每个态占波矢空间体积$(2\pi)^3/V$. 若自旋简并度$g=2s+1$, 球壳计数并换变量得
\[
D(\epsilon)=\frac{gV}{4\pi^2}\left(\frac{2m}{\hbar^2}\right)^{3/2}\sqrt\epsilon.
\]
硬壁正弦态只用第一卦限, 格距却缩小一半, 八分之一卦限与八倍格点密度相抵, 给同一体相计数. 零温每态至多占一粒子, 最低态从内到外填成Fermi球. 对自旋$1/2$粒子,
\[
k_F=(3\pi^2n)^{1/3},\quad \epsilon_F=\frac{\hbar^2k_F^2}{2m},\quad
U_0=\int_0^{\epsilon_F}\epsilon D(\epsilon)\dd\epsilon=\frac35N\epsilon_F.
\]
$\epsilon_F$是零温化学势, $T_F=\epsilon_F/\kb$是简并温度尺度, 并非实际温度. $T=0$时粒子仍有量级$v_F=\hbar k_F/m$的速度, Pauli约束迫使多数粒子离开最低单粒子态.

固定$N$时$U_0\propto V^{-2/3}$, 所以$P_0=-(\partial U_0/\partial V)_S=2n\epsilon_F/5$, $\kappa_T=3/(5P_0)$. 零温压力由状态计数产生, 不依赖经典热运动. 自由电子模型对简单金属能给出合理压缩率量级, 但晶格、带结构和相互作用会改变定量结果.

\originalsection{低温展开与Pauli顺磁性}
有限温下粒子数方程$N=\int D(\epsilon)f(\epsilon,\mu,T)\dd\epsilon$同时决定$\mu(T)$; 求$C_V$时必须追踪这个变化. 对平滑函数$g$且$\mu\gg\kb T$, 将Fermi函数与零温阶跃之差在$\epsilon=\mu$附近展开. 奇偶对称使常数项和二阶导数项消失, 一阶导数项带积分
\[
2\int_0^\infty\frac{x}{\e^x+1}\dd x
=2\sum_{r\ge1}\frac{(-1)^{r-1}}{r^2}=\frac{\pi^2}{6}.
\]
于是Sommerfeld展开为$\int_0^\infty g f\dd\epsilon=\int_0^\mu g\dd\epsilon+(\pi^2/6)(\kb T)^2g'(\mu)+O(T^4)$. 远离谱边缘且$g$足够光滑保证余项控制; 临近带边或强奇点时不能直接用它.

先取$g=D\propto\sqrt\epsilon$守恒$N$, 再取$g=\epsilon D$, 得
\[
\mu=\epsilon_F\left[1-\frac{\pi^2}{12}(T/T_F)^2+O(T^4/T_F^4)\right],
\quad C_V=\frac{\pi^2}{2}N\kb\frac T{T_F}+O(T^3).
\]
热容小的物理原因是只有Fermi面附近宽度$\kb T$中的粒子能改变占据, 比例约$T/T_F$, 不是每个电子都贡献$3\kb/2$.

弱场将两自旋能量移动$\mp\mu_eB$. 零温上下自旋粒子数差为$\mu_eBD(\epsilon_F)$, 其中$D$包括两自旋. 磁矩$M=\mu_e^2BD(\epsilon_F)$, 磁化率在$T\ll T_F$近似与温度无关. 若按体积磁化或对$H$求导, 还需相应$V$或SI的$\mu_0$因子. 这个Pauli顺磁规律不同于局域独立磁矩的Curie律, 因为只有Fermi面附近态参与重新分配.

\originalsection{能带、准粒子与半导体}
金属在化学势处有可用态; 绝缘体与本征半导体在价带顶和导带底之间有能隙. 带中电子不是完全自由粒子, 在极值附近可用有效质量$m^*$的抛物色散, 并据此改写态密度. 相互作用系统的低能激发还包括穿着相互作用云的准粒子以及集体模式: 晶格声子、电子等离子体振荡、铁磁自旋波、表面ripplon等. 它们的数目与色散不同, 不能一律套自由电子态密度.

取价带顶能量0, 导带底$E_g$, 本征材料有$N_e=N_h$. 在$E_g\gg\kb T$的非简并近似下,
\[
\begin{aligned}
 n_e&=2\left(\frac{m_e^*\kb T}{2\pi\hbar^2}\right)^{3/2}\e^{-(E_g-\mu)/(\kb T)},\\
 n_h&=2\left(\frac{m_h^*\kb T}{2\pi\hbar^2}\right)^{3/2}\e^{-\mu/(\kb T)}.
\end{aligned}
\]
它们来自带边$\sqrt\epsilon$态密度的Gaussian积分. 相等给$\mu=E_g/2+(3/4)\kb T\ln(m_h^*/m_e^*)$. 相同有效质量时化学势在隙中点; 不同质量时热态密度差使它偏移. 掺杂加入电中性约束和杂质态, 不再用$N_e=N_h$代替全部粒子数方程.
